% =====================================================================
%  TCET CODING PLATFORM -- Official Technical Documentation
%  MODULE 2 of 7 : System Architecture & Technology Stack
%  STANDALONE FILE.  Compile: pdflatex Module2_Architecture.tex  (twice)
% =====================================================================
\documentclass[12pt,a4paper,oneside]{report}
\usepackage[T1]{fontenc}
\usepackage[utf8]{inputenc}
\usepackage[a4paper,margin=1in]{geometry}
\usepackage{mathptmx}\usepackage{microtype}\usepackage{graphicx}\usepackage{xcolor}
\usepackage{tabularx}\usepackage{booktabs}\usepackage{longtable}\usepackage{xltabular}
\usepackage{array}\usepackage{enumitem}\usepackage{listings}\usepackage{titlesec}
\usepackage{fancyhdr}\usepackage{parskip}\usepackage[hidelinks]{hyperref}
\definecolor{tcetblue}{RGB}{0,0,0}\definecolor{tcetaccent}{RGB}{0,0,0}
\definecolor{codebg}{RGB}{245,246,248}\definecolor{codeframe}{RGB}{215,219,225}
\definecolor{codekw}{RGB}{0,0,0}\definecolor{codecomment}{RGB}{0,0,0}
\definecolor{codestr}{RGB}{0,0,0}\definecolor{boxframe}{RGB}{0,0,0}
\hypersetup{colorlinks=true,linkcolor=tcetblue,urlcolor=tcetaccent,citecolor=tcetblue,
  pdftitle={TCET Coding Platform -- Module 2: Architecture},pdfauthor={Centre of Excellence, TCET}}
\lstdefinestyle{tcp}{backgroundcolor=\color{codebg},basicstyle=\ttfamily\footnotesize,
  keywordstyle=\color{codekw}\bfseries,commentstyle=\color{codecomment}\itshape,
  stringstyle=\color{codestr},numberstyle=\tiny\color{gray},frame=single,
  rulecolor=\color{codeframe},breaklines=true,showstringspaces=false,tabsize=2,
  columns=fullflexible,keepspaces=true,xleftmargin=0.6em,framexleftmargin=0.6em}
\lstset{style=tcp}
\newcommand{\code}[1]{\texttt{#1}}
\newenvironment{callout}[1]{\par\vspace{4pt}\noindent
  \begin{tabular}{@{}p{0.014\linewidth}@{\hspace{0.6em}}p{0.93\linewidth}@{}}
  \textcolor{boxframe}{\rule{2pt}{\baselineskip}} & \textbf{\textcolor{tcetblue}{#1}}\\[2pt] & }{\end{tabular}\par\vspace{6pt}}
\titleformat{\chapter}[display]{\bfseries\fontsize{18}{22}\selectfont}{\bfseries\fontsize{14}{17}\selectfont Chapter \thechapter}{0.2em}{}
\titlespacing*{\chapter}{0pt}{6pt}{22pt}
\titleformat{\section}{\bfseries\fontsize{16}{19}\selectfont}{\thesection}{0.6em}{}
\titleformat{\subsection}{\bfseries\fontsize{14}{17}\selectfont}{\thesubsection}{0.5em}{}
\pagestyle{fancy}\fancyhf{}\renewcommand{\headrulewidth}{0.4pt}
\fancyhead[L]{\small\itshape TCET Coding Platform}\fancyhead[R]{\small\itshape Module 2 \textbullet\ Architecture}
\fancyfoot[C]{\thepage}
\fancypagestyle{plain}{\fancyhf{}\fancyfoot[C]{\thepage}\renewcommand{\headrulewidth}{0pt}}
\newcommand{\note}[1]{\par\vspace{3pt}\noindent{\color{tcetaccent}\textbf{Note.}}~#1\par\vspace{3pt}}
\begin{document}
\begin{titlepage}\centering
\vspace*{1.4cm}{\color{tcetaccent}\rule{\linewidth}{1.2pt}}\\[0.5cm]
{\LARGE\bfseries TCET Coding Platform\par}
\vspace{0.4cm}{\Large Official Technical Documentation\par}{\large Draft~1 \textemdash\ Testing Phase\par}
{\color{tcetaccent}\rule{\linewidth}{1.2pt}}\\[1.4cm]
{\Huge\bfseries\color{tcetblue} Module 2\par}\vspace{0.3cm}
{\Huge System Architecture \&\\[0.2cm] Technology Stack\par}\vspace{2.2cm}
{\large\bfseries Powered by\par}{\Large\color{tcetblue} Centre of Excellence (CoE), TCET\par}\vfill
\begin{tabular}{rl}\textbf{Module} & 2 of 7\\[3pt]
\textbf{Scope} & Architecture, topology, request lifecycle, technology\\[3pt]
\textbf{Status} & Draft~1 (living document until Stage~2)\\[3pt]
\textbf{Audience} & Officials, Architects, Engineering\\\end{tabular}
\vspace{0.8cm}{\small\itshape \today\par}\end{titlepage}
\pagenumbering{roman}\tableofcontents\clearpage\pagenumbering{arabic}

% =====================================================================
\chapter{Architectural Overview}
% =====================================================================
\section{Purpose and Scope of this Module}
This module specifies the architecture of the TCET Coding Platform:
its structural decomposition, the way requests flow through it, the topology of
the running processes, and the technologies chosen at each layer together with
the reasoning behind them. It is the structural companion to Module~1's
functional overview and the foundation on which the security, execution,
contest, proctoring, and operations modules build.

\section{Architectural Style}
The platform is a full-stack web application built in a \textbf{layered,
service-oriented} style with an \textbf{asynchronous execution subsystem}. Four
architectural commitments shape everything that follows:

\begin{description}[leftmargin=1.6em,style=nextline]
  \item[Separation of request handling from code execution] User requests are
        served synchronously and quickly; the expensive, untrusted work of
        compiling and running submitted code is pushed onto a queue and handled
        by background workers. This single decision is what lets the platform
        absorb mass-submission bursts without collapsing.
  \item[A hard trust boundary at the edge] Authentication is performed upstream
        by the CoE Single Sign-On; the backend trusts a reverse proxy to forward
        a verified identity and refuses traffic that does not arrive through it.
  \item[Layering with dependency inversion] Each domain is organized into
        routes, controllers, validators, services, repositories, and models.
        Business logic depends on interfaces, never on concrete infrastructure,
        so storage and execution engines are swappable and independently
        testable.
  \item[Deterministic, self-healing behaviour] Time-dependent logic uses an
        injected clock; finalization and scoring are idempotent; and background
        schedulers recover stranded work, so the system degrades gracefully.
\end{description}

\begin{callout}{The central idea}
A student's click on ``Submit'' does not wait for a compiler. It records a
submission, drops a job on a queue, and returns immediately. Somewhere else, a
worker picks up that job, runs the code inside a sandbox, and writes back the
verdict. Everything about the platform's ability to scale follows from this
separation.
\end{callout}

\section{Quality Attributes Addressed}
\begin{longtable}{@{}p{0.26\linewidth} p{0.66\linewidth}@{}}
\toprule
\textbf{Attribute} & \textbf{How the architecture addresses it}\\
\midrule
\endhead
Scalability & Queue-based buffering; horizontally scalable stateless workers.\\
Performance & Asynchronous submission keeps request latency independent of judging.\\
Security & Edge trust boundary; layered hardening; sandboxed execution.\\
Reliability & Idempotent finalization; stale recovery; process supervision.\\
Testability & Interface-based layers; injected clock; in-memory and stub doubles.\\
Maintainability & Uniform per-module layering; a single composition root.\\
Portability & Pluggable execution provider and repository implementations.\\
\bottomrule
\end{longtable}

% =====================================================================
\chapter{System Context \& Container View}
% =====================================================================
This chapter presents the architecture at two zoom levels: the \emph{context}
(the platform and the external systems it interacts with) and the
\emph{containers} (the major deployable/runtime units inside the platform).

\section{Context View}
\begin{lstlisting}[language={}]
                      +-------------------------+
    Student / Faculty |   Web browser (SPA)     |
        (human)  <--->|   React single-page app |
                      +-----------+-------------+
                                  | HTTPS
                                  v
                      +-------------------------+       +------------------+
                      | Reverse proxy / edge    |<----->| CoE SSO          |
                      | (Cloudflare Tunnel)     | auth  | (identity issuer)|
                      +-----------+-------------+       +------------------+
                                  | trusted source IP + signed identity
                                  v
                      +-------------------------+
                      |   TCET Coding Platform  |
                      |   (this system)         |
                      +-----------+-------------+
                        |         |          |
                        v         v          v
                   +--------+ +--------+ +-----------+
                   | MongoDB| | Redis  | | Judge0    |
                   | (data) | | (queue)| | (sandbox) |
                   +--------+ +--------+ +-----------+
\end{lstlisting}

\begin{longtable}{@{}p{0.24\linewidth} p{0.68\linewidth}@{}}
\toprule
\textbf{External entity} & \textbf{Relationship to the platform}\\
\midrule
\endhead
Web browser (SPA) & The user's client; renders the app and enforces client-side proctoring.\\
Reverse proxy / edge & Terminates public routing; forwards authenticated requests with identity headers.\\
CoE SSO & Authenticates users upstream; the sole source of identity.\\
MongoDB & System of record for all persistent state.\\
Redis & Backing store for the submission queue.\\
Judge0 & Sandboxed multi-language execution engine.\\
\bottomrule
\end{longtable}

\section{Container View}
Within the platform boundary, the following runtime containers cooperate:

\begin{longtable}{@{}p{0.24\linewidth} p{0.68\linewidth}@{}}
\toprule
\textbf{Container} & \textbf{Responsibility}\\
\midrule
\endhead
Frontend SPA & Serves and runs the React application (routing, editor, contest UI).\\
Backend API & The Express service: validation, authorization, business logic, persistence, enqueueing.\\
Submission worker & A background process that consumes the queue and drives judging and finalization.\\
Background schedulers & In-process timers (inside the API) for stale-submission recovery and expired-attempt finalization.\\
Datastore & MongoDB collections holding users, problems, submissions, contests, attempts, and more.\\
Queue & Redis-backed BullMQ queue of pending submissions.\\
Execution engine & Judge0 running the \code{isolate} sandbox.\\
\bottomrule
\end{longtable}

\note{The submission worker may run \emph{embedded} inside the API process (a
configuration switch) or as a \emph{separate} process. Production uses a separate
worker so judging load cannot affect request latency; see Chapter~7 and
Module~7.}

% =====================================================================
\chapter{Logical Components}
% =====================================================================
Zooming into the backend container, the platform is composed of the following
logical components. Each is described by its responsibility and its collaborators.

\section{Component Catalogue}
\begin{longtable}{@{}p{0.26\linewidth} p{0.66\linewidth}@{}}
\toprule
\textbf{Component} & \textbf{Responsibility \& collaborators}\\
\midrule
\endhead
HTTP application & Wires global middleware (security headers, CORS, guards, rate limiting) and mounts the domain routers.\\
Auth middleware & Verifies the trusted-proxy source and the forwarded identity; establishes the request user.\\
Domain routers & Bind endpoints to controllers and apply role/rate-limit middleware per route.\\
Controllers & Translate HTTP requests into service calls and shape responses.\\
Validators & Zod schemas that validate and type every inbound payload.\\
Services & The seat of business logic: users, problems, submissions, contests, leaderboard.\\
Repositories & The only components that talk to MongoDB; expose interfaces to services.\\
Domain models & Pure types and transformation/scoring functions with no I/O.\\
Execution provider & Abstraction over Judge0 (or a stub) for running code.\\
Submission queue & Abstraction over BullMQ for enqueueing submissions.\\
Schedulers & Timers that recover stale submissions and finalize expired attempts.\\
\bottomrule
\end{longtable}

\section{Component Interaction (submission path)}
\begin{lstlisting}[language={}]
Controller ---> SubmissionService.createSubmission()
                    |  validate problem visibility
                    |  ensure user record
                    |  persist submission (QUEUED)
                    |  SubmissionQueue.enqueue(id)  ----> Redis
                    v
                (returns receipt to client)

Worker ------> SubmissionService.processQueuedSubmission(id)
                    |  mark RUNNING
                    |  ExecutionProvider.executeSubmission()  ----> Judge0
                    |  finalizeSubmission(verdict)
                    |     update user + problem aggregates
                    |     update leaderboard
                    v
                (verdict persisted; client polling observes it)
\end{lstlisting}

\section{Dependency Rule}
Dependencies point inward: controllers depend on services; services depend on
repository, queue, and execution-provider \emph{interfaces}; models depend on
nothing. No inner layer imports an outer layer, and no business logic imports the
database driver. This is what makes the domain logic unit-testable with
in-memory doubles.

% =====================================================================
\chapter{Request Lifecycle \& Middleware Pipeline}
% =====================================================================
Every API request passes through an ordered pipeline before it reaches domain
logic. The order is deliberate: cheap, security-critical checks run first.

\section{The Pipeline in Order}
\begin{longtable}{@{}p{0.05\linewidth} p{0.28\linewidth} p{0.55\linewidth}@{}}
\toprule
\textbf{\#} & \textbf{Stage} & \textbf{Purpose}\\
\midrule
\endhead
1 & Security headers (Helmet) & Set safe response headers; disable \code{x-powered-by}.\\
2 & CORS & Admit only allowlisted origins; handle preflight.\\
3 & Cookie parsing & Make federated cookies available for identity.\\
4 & Body parsing (capped) & Parse JSON up to a 100\,kb limit.\\
5 & Global rate limiter & Throttle by identity or IP.\\
6 & Frontend-pathname guard & Validate the optional client pathname header.\\
7 & Mutation-origin guard & Require an allowlisted Origin/Referer on state-changing calls.\\
8 & Auth middleware (per router) & Verify trusted proxy + identity; attach the user.\\
9 & Role / profile guards (per route) & Enforce role and profile-completion requirements.\\
10 & Controller $\to$ service & Execute domain logic.\\
11 & Error handler & Convert errors into consistent HTTP responses.\\
\bottomrule
\end{longtable}

\section{Request Walkthrough (authenticated GET)}
\begin{lstlisting}[language={}]
GET /api/problems
  -> Helmet sets headers
  -> CORS: Origin allowlisted? yes
  -> rate limiter: under limit? yes
  -> pathname guard: header valid/absent? yes
  -> mutation guard: method is GET (safe) -> skip
  -> auth: source IP trusted? identity valid? status ACTIVE? -> attach user
  -> controller: ProblemService.listProblems(user, query)
  -> service: filter by visibility (published + department) and paginate
  -> 200 OK with the page of problems
\end{lstlisting}

\section{Error Handling Strategy}
Domain code throws typed application errors carrying an HTTP status and message.
A single terminal error handler converts these into consistent JSON responses,
so controllers never assemble error bodies themselves and unexpected exceptions
never leak stack traces to clients. A not-found handler returns a uniform 404 for
unmapped routes.

% =====================================================================
\chapter{Backend Layered Architecture}
% =====================================================================
The backend is organized so that every domain module has the same internal
shape. This uniformity is a maintainability feature: a reader who understands one
module understands them all.

\section{The Six Layers}
\begin{longtable}{@{}p{0.22\linewidth} p{0.70\linewidth}@{}}
\toprule
\textbf{Layer} & \textbf{Responsibility}\\
\midrule
\endhead
Routes & Declare endpoints; attach auth, role, and rate-limit middleware; delegate to controllers.\\
Controllers & Adapt HTTP to the service API; read validated input; write HTTP responses.\\
Validators & Define Zod schemas; validation produces both a runtime guarantee and a static type.\\
Services & Hold all business rules; orchestrate repositories, the queue, and the execution provider.\\
Repositories & Encapsulate persistence; map between documents and domain records; the only MongoDB callers.\\
Models & Pure domain types and functions (e.g. scoring), free of I/O and framework concerns.\\
\bottomrule
\end{longtable}

\section{Domain Modules}
\begin{longtable}{@{}p{0.20\linewidth} p{0.72\linewidth}@{}}
\toprule
\textbf{Module} & \textbf{Responsibility}\\
\midrule
\endhead
User & Identity synchronization, profile, and profile analytics.\\
Problem & Problem authoring, lifecycle, visibility, and aggregates.\\
Submission & The run/submit pipeline, judging orchestration, and finalization.\\
Contest & The full contest engine: attempts, scoring, standings, feedback.\\
Leaderboard & Denormalized ranking derived from user statistics.\\
\bottomrule
\end{longtable}

\section{Why Layering Pays Off Here}
\begin{itemize}[leftmargin=1.5em]
  \item \textbf{Testability.} Because services depend on repository interfaces,
        tests substitute in-memory repositories and a stub execution provider and
        exercise real business logic deterministically.
  \item \textbf{Swappability.} The datastore or execution engine can change
        behind its interface without touching business logic.
  \item \textbf{Locality of change.} A change to validation, HTTP shape, or a
        business rule lives in exactly one layer.
\end{itemize}

% =====================================================================
\chapter{Dependency Injection \& Composition}
% =====================================================================
\section{The Composition Root}
All concrete dependencies are assembled in one place \textemdash\ a composition
root that constructs the repositories, the execution provider, the submission
queue, and every service, wires them together, and returns a single
application-dependencies object that both the API and the worker consume.

\section{Overrides for Testing}
The composition root accepts overrides for the auth middleware, execution
provider, submission queue, repositories, and the clock. This is the mechanism
that makes the whole system testable without production infrastructure: a test
builds the application with in-memory repositories, a stub execution provider,
and a fixed clock, and then exercises the real routes and services end to end.

\begin{callout}{The injected clock}
Time is a dependency, not a global. Services obtain ``now'' from an injected
function. In production it returns the wall clock; in tests it returns a fixed
instant. This is why contest windows, personal deadlines, stale-recovery
thresholds, and finalization are all deterministically testable.
\end{callout}

\section{Startup Sequence}
\begin{lstlisting}[language={}]
1. Load & validate configuration (fail fast on invalid config).
2. Build application dependencies (repositories, services, provider, queue).
3. Create the Express app (mount middleware + routers).
4. Optionally create an embedded submission worker.
5. Start listening; run stale-submission recovery once.
6. Start the expired-attempt finalizer timer (unless disabled).
7. Register graceful-shutdown handlers.
\end{lstlisting}

% =====================================================================
\chapter{Process Topology}
% =====================================================================
\section{Runtime Processes}
\begin{longtable}{@{}p{0.24\linewidth} p{0.68\linewidth}@{}}
\toprule
\textbf{Process} & \textbf{Role}\\
\midrule
\endhead
API server & Serves the REST API; runs the background schedulers; may embed a worker.\\
Submission worker & Dedicated queue consumer; recommended as a separate process in production.\\
Frontend & Serves the built single-page application.\\
\bottomrule
\end{longtable}

\section{Embedded vs. Separate Worker}
The API can host a worker in-process for simple deployments, or delegate all
judging to one or more standalone worker processes. Production separates them so
that a surge of judging work \textemdash\ CPU-heavy and bursty \textemdash\
cannot degrade API responsiveness. The two topologies are shown below.

\begin{lstlisting}[language={}]
Simple (single process):        Production (separated):
+---------------------+         +-------------+   +-----------------+
| API + embedded      |         | API server  |   | Worker process  |
| worker              |         | (schedulers)|   | (concurrency N) |
+----------+----------+         +------+------+   +--------+--------+
           |                           |                   |
        Redis                        Redis  <------------- Redis
\end{lstlisting}

\section{Graceful Shutdown}
On receiving a termination signal the API stops accepting connections, clears the
finalizer timer, and closes the embedded worker (if any) before exiting, so
in-flight work is not abruptly severed.

% =====================================================================
\chapter{The Asynchronous Execution Subsystem}
% =====================================================================
\section{Why a Queue}
Compiling and running untrusted code is slow and resource-intensive relative to a
typical web request, and examination load is bursty. Executing on the request
path would tie up connections and collapse under a mass submit. A queue converts
a synchronous, unbounded demand into an asynchronous, bounded one: submissions
are accepted instantly and drained at a controlled rate.

\section{Queue Semantics}
\begin{longtable}{@{}p{0.28\linewidth} p{0.64\linewidth}@{}}
\toprule
\textbf{Property} & \textbf{Behaviour}\\
\midrule
\endhead
Job identity & Each submission is enqueued under its own identifier.\\
Retries & A bounded number of automatic attempts with exponential backoff.\\
Concurrency & Workers process several jobs in parallel, tunable and capped.\\
Completion retention & Completed jobs are retained briefly for observability.\\
Idempotency & Finalization is guarded so a re-delivered job cannot double-apply a verdict.\\
\bottomrule
\end{longtable}

\section{Backpressure \& Draining}
When submissions arrive faster than they can be judged, they accumulate in Redis
rather than failing. Throughput is governed by the number of worker processes and
their concurrency, which are matched to what the Judge0 stack can sustain. Adding
workers increases drain rate horizontally.

\section{Self-Healing}
\begin{itemize}[leftmargin=1.5em]
  \item \textbf{Stale-submission recovery} re-enqueues submissions left in a
        pending state past a staleness threshold, protecting against a worker
        that died mid-judging.
  \item \textbf{Expired-attempt finalization} finalizes contest attempts whose
        deadline has passed, so abandoned attempts are graded near their real
        deadline. Both are covered in Modules~4 and~5.
\end{itemize}

% =====================================================================
\chapter{Frontend Architecture}
% =====================================================================
\section{Application Shape}
The frontend is a Vite-built React single-page application in TypeScript. It is
organized into pages (per role), reusable components, hooks, a typed API client,
and small pure libraries for formatting and domain helpers.

\section{Routing \& Role Gating}
All routes are declared centrally. Authenticated routes are wrapped in a role
guard that admits only the permitted role(s), producing cleanly separated Student
and Faculty areas. A representative route map:

\begin{longtable}{@{}p{0.52\linewidth} p{0.36\linewidth}@{}}
\toprule
\textbf{Route} & \textbf{Role}\\
\midrule
\endhead
\code{/} & Public entry / login\\
\code{/complete-profile} & Student \& Faculty\\
\code{/student/dashboard}, \code{/student/problems} & Student\\
\code{/student/problems/:id} & Student\\
\code{/student/contests}, \code{/student/contests/:id} & Student\\
\code{/student/contests/:id/questions/:qid} & Student\\
\code{/student/contests/:id/feedback} & Student\\
\code{/student/leaderboard}, \code{/student/profile} & Student\\
\code{/faculty/dashboard}, \code{/faculty/problems} & Faculty\\
\code{/faculty/create-problem}, \code{/faculty/create-contest} & Faculty\\
\code{/faculty/contests}, \code{/faculty/contests/:id} & Faculty\\
\code{/faculty/submissions}, \code{/faculty/leaderboard} & Faculty\\
\bottomrule
\end{longtable}

\section{Server State \& Data Fetching}
Server state is managed with a query library that caches responses, deduplicates
requests, and refetches in the background, keeping the UI responsive and
consistent. A single typed API client centralizes every backend call and maps
responses into view models, so components never construct URLs or parse raw
responses directly.

\section{UI System \& Editor}
The interface uses a utility-first CSS framework with an accessible component
library, supporting light and dark themes. The in-browser Monaco editor provides
the coding experience with per-language syntax support. Dedicated contest
components enforce the examination environment (timer, question navigation,
watermark, screen guard, lock overlay) and integrate with the proctoring hook
described in Module~6.

\section{Client-Side Resilience}
The client polls for asynchronous submission results, shows optimistic and
loading states, and guards contest navigation. Route-level error boundaries
contain rendering failures so a single faulty view does not blank the whole app.

% =====================================================================
\chapter{Data Flow \& Sequence Walkthroughs}
% =====================================================================
This chapter traces three representative flows end to end.

\section{Sequence: Practice Submission}
\begin{lstlisting}[language={}]
Student        Frontend        API            Queue     Worker     Judge0
  |   write code   |             |              |          |          |
  |--------------->|             |              |          |          |
  |   Submit       |             |              |          |          |
  |--------------->| POST submit |              |          |          |
  |               |------------>| persist QUEUED|         |          |
  |               |             |---- enqueue ->|         |          |
  |               |<- receipt --|              |          |          |
  |   poll status  |             |              |          |          |
  |--------------->|-- GET ----->|              |          |          |
  |               |             |              |-- job -->|          |
  |               |             |              |          |-- run -->|
  |               |             |              |          |<- verdict|
  |               |             |<-- finalize (persist verdict) -----|
  |<- ACCEPTED ----|<- status ---|              |          |          |
\end{lstlisting}

\section{Sequence: Authenticated Access}
\begin{lstlisting}[language={}]
Browser -> Edge/SSO: request (with session)
Edge/SSO -> API:     forward + x-coe-* identity  (from trusted source IP)
API auth:            source trusted? identity valid? status ACTIVE?
API:                 sync user record; attach req.user
API -> Browser:      role-appropriate response
\end{lstlisting}

\section{Sequence: Contest Attempt (high level)}
\begin{lstlisting}[language={}]
register -> start attempt (freeze deadline) -> answer/save/submit (no scoring)
   -> end (submit | deadline | violation) -> auto-submit drafts
   -> faculty publish (authoritative re-grade) -> feedback gate -> report/standings
\end{lstlisting}

% =====================================================================
\chapter{Technology Stack Deep-Dive}
% =====================================================================
\section{Backend Technologies}
\begin{longtable}{@{}p{0.22\linewidth} p{0.70\linewidth}@{}}
\toprule
\textbf{Technology} & \textbf{Role \& rationale}\\
\midrule
\endhead
Node.js & Event-driven runtime well suited to I/O-bound API and queue work.\\
Express~5 & Minimal, mature HTTP framework; explicit middleware model.\\
TypeScript & Static typing across the whole domain reduces defects and documents contracts.\\
Zod & Runtime validation that also yields static types; one schema secures and types a payload.\\
Helmet & Secure HTTP headers with minimal configuration.\\
Rate limiter & Protects the API from floods; keyed by identity or IP.\\
MongoDB driver & Direct, efficient access to the document store.\\
BullMQ & Robust Redis-backed job queue with retries and concurrency.\\
\bottomrule
\end{longtable}

\section{Frontend Technologies}
\begin{longtable}{@{}p{0.22\linewidth} p{0.70\linewidth}@{}}
\toprule
\textbf{Technology} & \textbf{Role \& rationale}\\
\midrule
\endhead
React~18 & Component model for a complex, stateful UI.\\
Vite~5 & Fast build and dev tooling.\\
TypeScript & End-to-end typing shared with API view models.\\
React Router & Declarative routing and role gating.\\
TanStack Query & Server-state caching, deduplication, and background refetch.\\
Tailwind + component library & Consistent, accessible, themeable design system.\\
Monaco editor & Professional in-browser code editing.\\
\bottomrule
\end{longtable}

\section{Data \& Execution Technologies}
\begin{longtable}{@{}p{0.22\linewidth} p{0.70\linewidth}@{}}
\toprule
\textbf{Technology} & \textbf{Role \& rationale}\\
\midrule
\endhead
MongoDB & Flexible document model fits nested, snapshot-heavy records (attempts embed per-question state); server-side JSON Schema validation adds a safety net.\\
Redis & Low-latency backing store for the queue and buffering.\\
Judge0 + \code{isolate} & Purpose-built sandboxed execution for many languages; runs untrusted code safely.\\
\bottomrule
\end{longtable}

% =====================================================================
\chapter{Cross-Cutting Concerns}
% =====================================================================
\section{Configuration}
All configuration is centralized and validated at startup against a schema, so
the process fails fast rather than starting in a partially-configured or insecure
state (for example, an over-short signing secret is rejected before the server
listens). The full variable reference is in Module~7.

\section{Logging \& Observability}
Security-relevant events (authentication rejections) and background-scheduler
outcomes (recovered submissions, finalized attempts) are logged to the process
console for aggregation. Health endpoints are restricted to trusted internal
sources. Richer structured logging and metrics are identified as future work in
Module~7.

\section{Error Handling}
A uniform error model (typed errors with HTTP status) and a single terminal
handler ensure consistent responses and prevent leakage of internal detail.

\section{Resilience}
Idempotent finalization, stale-submission recovery, and the expired-attempt
finalizer combine to make the platform tolerant of worker failure, queue
backpressure, and student disconnection.

% =====================================================================
\chapter{Network, Ports \& Environments}
% =====================================================================
\section{Ports}
\begin{longtable}{@{}p{0.24\linewidth} p{0.16\linewidth} p{0.50\linewidth}@{}}
\toprule
\textbf{Service} & \textbf{Port} & \textbf{Notes}\\
\midrule
\endhead
Frontend & 5173 & Dev / preview server.\\
Backend API & 3000 / 3001 & Configurable; local dev often 3001.\\
Mock CoE SSO & 4000 & Local development identity stub only.\\
Judge0 & 2358 & Local sandboxed execution API.\\
Redis & 6379 & Submission queue.\\
MongoDB & 27017 & Datastore.\\
\bottomrule
\end{longtable}

\section{Environments}
\begin{longtable}{@{}p{0.20\linewidth} p{0.72\linewidth}@{}}
\toprule
\textbf{Environment} & \textbf{Characteristics}\\
\midrule
\endhead
Development & Local services; a mock SSO; a stub or local Judge0; embedded worker acceptable.\\
Test & In-memory repositories, stub execution, fixed clock; deterministic.\\
Production & Reverse proxy + real SSO; MongoDB, Redis, and Judge0; separated worker; PM2 supervision.\\
\bottomrule
\end{longtable}

% =====================================================================
\chapter{Architecture Decision Records}
% =====================================================================
The following condensed ADRs record the principal architectural decisions, the
context that motivated them, and their consequences.

\section{ADR-1: Asynchronous submission via a queue}
\textbf{Context:} mass, bursty submission load; slow, untrusted execution.
\textbf{Decision:} accept submissions synchronously and judge them on a queue.
\textbf{Consequences:} request latency is decoupled from judging; the system
scales by adding workers; clients poll for results; finalization must be
idempotent.

\section{ADR-2: Trusted reverse-proxy identity}
\textbf{Context:} the institution already operates a centralized SSO.
\textbf{Decision:} delegate authentication upstream and trust forwarded identity
from an allowlisted proxy. \textbf{Consequences:} no credential storage; a hard
requirement never to expose the backend directly; a source-IP allowlist becomes
security-critical.

\section{ADR-3: Layered architecture with dependency inversion}
\textbf{Context:} a need for testable, maintainable domain logic.
\textbf{Decision:} uniform routes/controller/validator/service/repository/model
layering with interface-based dependencies and an injected clock.
\textbf{Consequences:} high testability and swappability; a small amount of
boilerplate per module.

\section{ADR-4: MongoDB with snapshots}
\textbf{Context:} nested, history-sensitive records. \textbf{Decision:} a
document store with denormalized snapshots and server-side schema validation.
\textbf{Consequences:} cheap reads and stable history; aggregates must be
recomputed on finalization.

\section{ADR-5: Pluggable execution provider}
\textbf{Context:} the need to test without Judge0 and to support local/cloud
engines. \textbf{Decision:} an execution-provider interface with Judge0 and stub
implementations. \textbf{Consequences:} deterministic tests; portability across
execution backends.

\vspace{0.6cm}
\begin{center}
\textit{--- End of Module 2. Continued in Module 3: Authentication, Authorization, Security \& Data Persistence. ---}
\end{center}
\end{document}
